% ch2 B.2 元胞法 (原书页 43–52 / PDF p058–p067)
% 衔接说明：p058 页首即本小节标题，其上无前节残句，本文件自标题起译。

\subsection{元胞法：结合能的定量计算}

这是一种完全不同的方法；在许多方面，它都已证明是计算金属内聚能（cohesive energy）迄今为止最为准确的方法。元胞法（cellular method）由 Wigner 和 Seitz 首先采用，可惜由于种种已知的原因，它似乎只容易应用于单价碱金属。

元胞法在几乎每一个可能的方面都与基于平面波的方法形成对照。在近自由电子模型中，以及我们将会看到的、在 O.P.W.（正交化平面波）方法中，波函数从一开始就被选为恰当的能带函数，满足晶格周期性所强加的边界条件；随后的功夫则在于用微扰或其他方法足够精确地处理离子实（ion cores）的势对这些函数产生的效应。元胞法却是从相反的尝试出发：先在离子实区域求出尽可能精确的势和尽可能精确的解，然后再去近似地满足由周期性所强制的边界条件。看来，在计算能量时这是更为精确的途径，但它只是对固体 s 带中最低的波函数才是如此。

具体做法是把实空间格子划分成元胞，方式与在倒格子中构成布里渊区的做法完全相同：对指向近邻的矢量作平分。对于碱金属通常被假定具有的 b.c.c. 格子，这样得到的是由 $\tfrac{1}{2}\,\tfrac{1}{2}\,\tfrac{1}{2}$ 平面构成的正八面体，其角被 $\tfrac{1}{2}\,00$ 平面切去——总共 14 个面——正是我们已经熟悉的那种多面体。元胞法的想法是：这个多面体其实同球相去不远，而且在发生偏离的外部区域，势非常弱、波函数非常平滑。于是就作出了元胞严格为球的近似。这使得对低态数值积分波动方程的问题得到如下若干简化：

1. 格子的其余部分也不妨同样假定是球对称的；由于碱金属中原子间的距离使得离子实几乎不发生交叠，来自其他元胞的势纯粹是静电势，而在这个假定下它完全消失，因为格子整体是中性的。

2. 再者，原子芯是闭合壳层，因而是球形的；结果我们在元胞内要求解的乃是一个球对称的——即完全可分离的——波动方程。

3. 对于 s 带的最低态——它当然 $k=0$，因而其波函数是简单周期的，即 $u_0(r)$——边界条件取如下简单形式：斜率 $\mathrm{d}\psi/\mathrm{d}r\,\big|_{r=R_{\mathrm{cell}}}=0$。这是因为波函数必须水平地趋近元胞边界，因为根据反射对称性，它必须关于边界为偶。最低的原子态则是 $\mathrm{d}\psi/\mathrm{d}r\,\big|_{r=\infty}=0$ 的那个球对称态。碱金属内聚、而且实际上也就是所谓“金属键”的根本物理基础在于：$\mathrm{d}\psi/\mathrm{d}r\,\big|_{r=R_c}=0$ 的态，其能量比自由原子中态的能量低相当多。Wigner 称之为“边界修正”（boundary correction），并在他的《Seitzschrift》文章 (25) 中对多种金属作了估计。

从相当粗略的考虑出发，我们立刻就能看出必然如此。球对称性使我们能够把波动方程 $-\nabla^2\Psi+V(r)\Psi=E\Psi$ 分离成角向方程和径向方程，只须假定 $\Psi=[U_L(r)/r]\,Y_{LM}(\theta,\phi)$；于是 $U_L$ 满足的方程是
\begin{equation*}
\frac{d^2 U_L}{dr^2}+\left[E-V(r)-\frac{L(L+1)}{r^2}\right]U_L=0
\end{equation*}
这里的单位是原子单位：$r$ 以玻尔半径计，所有能量以 rydberg 计。在碱金属原子的芯区之外，$V=-2Z/r$，在单价原子的情形即 $-2/r$；芯内的 $V$ 则更为复杂，而且实际上并不是一个局域算符，还包含与芯电子的交换。尽管如此，为了定性的目的，我们可以把它当作一个纯势 $-2Z(r)/r$ 来处理（图 8）。

%FIG{8}: 图 8

我们感兴趣的是低的 s 态（$L=0$），于是在 $r=0$ 处的边界条件是 $\Psi$ 有限，即 $U=0$ 且线性增大。让我们取某个很低的能量 $E_1$，在想象中从原点处这一边界条件出发把方程向外积分。在到达 $R_t(E_1)$ 处的“经典转折点”（classical turning point）之前，$E_1-V>0$，$\mathrm{d}^2U/\mathrm{d}r^2$ 与 $U$ 反号，曲率朝向轴；然而此后，$U>0$，$\mathrm{d}^2U/\mathrm{d}r^2>0$，曲率背离轴：$U$ 发散。现在让我们依次画出能量递增的 $E_1$、$E_2$、$E_3$ 以及 $E_f$——实际的自由原子本征值——的情形。我们看到，自由原子本征值高于使波函数在 $R_{\mathrm{cell}}$ 处恰好变平的那个能量；在此之后，波函数便转而增大并发散（图 9）。由此得到的动能收益可以相当可观。以钠为例，它是 3 1/2 伏特（译注：原文单位为 volts，此处即指电子伏）；在 Li 中更高，在重碱金属 K、Rb 和 Cs 中则更低。

%FIG{9}: 图 9

从这幅图出发，还可以对 Kuhn、Van Vleck、Brooks (21) 和 Ham (18) 为计算这一能量而发展的既精确又巧妙的方法获得相当好的了解。在两幅图中，我都标明了一个对所有碱金属都成立的事实：芯轨道（Li 中的 1s、Na 中的 2s2p，等等）的尺度比原子元胞小得多。这意味着在径向积分中存在一段相当可观的区域，在其中我们确切知道必须使用的势，即 $V=-2/r$。当然，芯内的势是未知的；在较早的计算中，对这一区域的做法是先取一个 Hartree 或 Hartree-Fock 势，然后修改它，使之拟合观测到的光谱项值序列——对碱金属原子，这些项值已为人们以极高的精度所知。

Kuhn 和 Van Vleck 指出，事实上也许永远不必穿过这一区域去积分径向方程。既然反正要使用一组取自光学光谱的经验能级来调整势，难道就不能有办法直接从光学光谱本身提取必要的信息吗？

Kuhn-Van Vleck 方法在原理上比 Brooks 和 Ham 最终完成的更精确、更完备的技术要简单。他们的做法是：把波函数看作从 $R=0$ 处 $U=0$ 出发、一直积分到 $R_{\mathrm{core}}$ 点的结果。在该点上，要把积分继续下去所需的全部信息就是对数导数（logarithmic derivative）$(1/U)(\mathrm{d}U/\mathrm{d}r)\,\big|_{r=R_{\mathrm{core}}}$。这个数将是初始能量 $E$ 的某种函数；如我们所见，对大的负 $E$，它从某个高值出发，下降穿过零，最后变得像一条 $\cot^{-1}$ 曲线。从 $R_{\mathrm{core}}$ 点往外，我们确切地知道势具有 $(-2/r)$ 的形式，而 Wannier 等人已经求出了任意能量下这两个解的解析形式。于是，我们可以把这些解拟合到

%FIG{10}: 图 10

任意给定的 $(1/U)(\mathrm{d}U/\mathrm{d}r)$ 值，并确定解在任意能量下的性质。特别地，我们能解析地定出那些可以导致 $\mathrm{d}U/\mathrm{d}r\,\big|_{r\to\infty}=0$ 的 $(1/U)(\mathrm{d}U/\mathrm{d}r)\,\big|_{R_{\mathrm{core}}}$ 值；它们就是观测到的光谱项能量处的值。随后我们就可以把这些值描绘出来；图 10 中的那些 $x$ 就是一组假想的值。在大致知道曲线应有的形状之后，我们可以试着过这些给定的点画一条光滑曲线，并假定它就是

正确的 $(1/U)(\mathrm{d}U/\mathrm{d}r)\,\big|_{r=R_c}(E)$。

我们还能解析地知道——这同样是因为我们掌握着库仑波方程解的精确解析行为——$(1/U)(\mathrm{d}U/\mathrm{d}r)\,\big|_{R_{\mathrm{core}}}$ 取什么值才会导致 $(1/U)(\mathrm{d}U/\mathrm{d}r)\,\big|_{R_{\mathrm{cell}}}=0$。把这个值记作 $(1/U)(\mathrm{d}U/\mathrm{d}r)\,\big|_0$。然后只须把我们的光滑曲线内插到这个值，便得到最低本征值的数值 $E_0$。

事实上，由于对数导数曲线的形状相当难以驾驭，它并\emph{不}是应当描绘的正确量；正如 Kuhn 本人首先指出的，能给出真正光滑外推的正确量是\emph{量子亏损}（quantum defect），它在自由原子本征值处正是碱金属光谱的 Ritz 内插公式中的量 $\delta_m$，
\begin{equation*}
E_m=-\frac{1}{(m-\delta_m)^2}
\end{equation*}
$\delta_m$ 与碰撞理论中出现的相移 $\delta$ 密切相关，不过库仑势的长程性质引入了一些除专家以外无人感兴趣的额外复杂因素。无论如何，人们发现——而且可以在数学上加以证明——总能找到一个量 $\delta$，它在观测到的光谱系各项之间光滑地内插，并且可用来满足元胞表面上任意希望的边界条件。在 Ham 的文章中，对五种碱金属各自的若干种密度，以三到四位有效数字给出了相应的 $k=0$ 能量。

这种方法的主要优点——这也可能是大多数内聚能计算（以别于实际能带形状的计算）用其他方法之所以不成功的最重要的原因之一——在于：价电子与芯电子相互作用中的多体效应，基本上被精确地包含在所使用的量子亏损经验值之中。甚至有可能在 $R_{\mathrm{core}}$ 之内根本不存在任何合理的局域势 $V(r)$ 能给出观测到的本征值谱；正如我们提到过的，即使在 H-F 中，正确的势也是非局域的。这看来实际上就是如下情形的主要原因：早期的 Wigner-Seitz 计算对 Li 和 Na 相当好，对 K 相当差，对 Rb 和 Cs 则根本行不通，而 Q.D.M.（量子亏损方法）对付重金属却毫无困难——在重金属中，由于芯更大、束缚更松，芯极化（core polarization）与交换都是很大的效应。

于是，$k=0$ 处那单独一条能级的定位，由 Q.D.-元胞法做得极为出色；这也正是它胜过所有其他方法的地方：至少有一条能级可以按绝对值准确定位。尽管如此，这还不可能是全部的战斗；我们必须首先计入所谓的“费米能”，即由于其余电子不得不进入带中较高的状态而出现的额外动能。接着，我们还必须对电子相互作用作一组相当大的修正；不过，这些修正之和结果证明相当小。

费米能是一个相当大的效应，从 Wigner 的定性论证 (20) 看得很清楚：该论证提示，一个 s 带被完全填满的物质将几乎是不结合的。从积分径向方程的观点看，带顶的状态必须是那种从原子到其近邻改变符号的状态，这意味着 $\psi$ 在 $R_{\mathrm{cell}}$ 处为零，而不是 $\mathrm{d}\psi/\mathrm{d}r$ 在 $R_c$ 处为零。$\psi=0$ 的波函数，其能量比自由原子波函数 $E_I$ 的能量所高出的量，大概与 $E_0$ 比 it 低的量大致相同。可以看到，在某种意义上，$\psi$ 在 $\infty$ 处趋于 0 大致是 $\psi\big|_{R_c}=0$ 与 $\mathrm{d}\psi/\mathrm{d}R\,\big|_{R_c}=0$ 之间的一个中间判据（图 11）。于是，对半满的带，我们或许会预期在费米能上损失掉 $k=0$ 态所赢得的 3 到 4 伏特能量的一半或更多，而事实证明确实如此。

%FIG{11}: 图 11

要计算费米能，就必须对 $k$ 空间中整个被占据区域的能带结构有相当准确的估计。结果表明，用元胞法时，对第一区边界附近的点、或者对波函数在元胞内不再呈 s 型的较高的带，这是一种困难而且常常不准确的计算。原因在于，元胞内波动方程的解正是角动量 $L=0,1,2$ 等等的 $s$、$p$、$d$ 等本征函数。……随着 $L$ 增大，计算或外推的精度迅速下降，因此，对 $L\geq 3$ 或 4 的波函数拟合规定的边界条件是非常困难的事情。这是由于有效排斥势能 $L(L+1)/r^2$ 的存在——在相应的光学能级中，它把电子挡在芯区之外；而固体问题中作边值拟合所需的波函数却确实穿透到芯区。

结果证明，遗憾的是，对相当大的 $k$，能带波函数包含非常可观的一些成分，必须用高角动量值的波函数才能描述。看到这一点的最简易途径是 Shockley 的“空晶格”（empty lattice）检验——我们可以把对称化平面波（它们是能带波函数的一个可能的集合）绕元胞中心按球谐函数展开，而且的确发现，对接近或越过区边界的那些 $k$ 值，它们含有相当可观的 $L\geq 3$ 或 4 的分量。尽管在把球谐函数拟合到正确边界条件这一数学问题上已经取得进展，但对碱金属以外的一切金属，其精度仍属可疑。

当然，对碱金属而言这个问题并不十分严重，因为很大一部分电子的 $k$ 值深处于第一区之内。这里使用的技巧是把能量 $E(k)$ 在原点附近按 $k$ 的幂级数展开；由对称性，第一项为零，第二项即所谓的有效质量（effective mass）项，
\begin{equation*}
E=E_0+\frac{\hbar^2 k^2}{2m^*}+\cdots
\end{equation*}
则可以用“$k\cdot p$”微扰论的一个计算求得。这种方法的原理是：写出并非由总波函数 $\psi_k(r)$、而是仅由其周期部分 $u_k$ 所满足的波动方程。我们从通常的波动方程
\begin{equation*}
\left[-\frac{\hbar^2}{2m}\nabla^2+V(r)\right]u_k(r)\,e^{ik\cdot r}=E_k\,u_k\,e^{ik\cdot r}
\end{equation*}
出发。显然，除 $\nabla^2$ 之外的一切都与 $e^{ik\cdot r}$ 对易，因此在完成如下的计算之后，指数因子便可以提出来：
\begin{align*}
\nabla^2\left(u_k\,e^{ik\cdot r}\right)&=\nabla\cdot\left(\nabla u_k+ik\,u_k\,e^{ik\cdot r}\right)\\
&=-k^2 u_k\,e^{ik\cdot r}+2ik\,e^{ik\cdot r}\cdot\nabla u_k+\left(\nabla^2 u_k\right)e^{ik\cdot r}
\end{align*}
于是
\begin{equation*}
-\frac{\hbar^2}{2m}\left(\nabla^2+2ik\cdot\nabla\right)u_k+V(r)\,u_k=\left(E_k-\frac{\hbar^2 k^2}{2m}\right)u_k
\end{equation*}
当我们把它与本征函数 $\psi_0=u_0$ 所满足的方程
\begin{equation*}
-\frac{\hbar^2}{2m}\nabla^2 u_0+V(r)\,u_0=E_0\,u_0
\end{equation*}
相比较时，我们看到，除了
\begin{equation*}
-i\,\frac{\hbar^2}{m}\,k\cdot\nabla u_k=\frac{\hbar}{m}\,k\cdot p\,u_k
\end{equation*}
这一项之外，它就是同一个方程，而且由于 $u_k$ 是周期的，还具有相同的边界条件。于是，对足够小的 $k$，我们可以把这一项当作微扰来处理，从而得到用 $k=0$ 处的解表出的 $k$ 处能量的表达式。事实上，只须稍加修改，这个表达式就可以在任意一般点 $k_0$ 处使用，以求出邻近点上的波函数与本征值。

在我们所讨论的碱金属这一特殊情形中，$k\cdot p$ 项含有梯度算符，因而把偶函数变为奇函数，反之亦然。这意味着对角矩阵元 $(u_{0n},\,(k\cdot p)\,u_{0n})$ 为零，所以一级修正为零。二级修正按标准微扰论为
\begin{equation*}
E_k-\left(\frac{\hbar^2 k^2}{2m}+E_0\right)=-\frac{\hbar^2 k^2}{m^2}\sum_n\frac{|(0n|p_x|00)|^2}{E_n(0)-E_0}
\end{equation*}
于是
\begin{equation*}
\frac{1}{m^*}=\frac{1}{m}\left[1-\frac{2}{m}\sum_n\frac{|(0n|p_x|00)|^2}{E_n(0)-E_0}\right]
\end{equation*}
乍看之下，有效质量似乎总应增大——能带变得更为平坦——因为右边的求和预期 $>0$，其中 $E_n>E_0$。对 Li 来说确实如此：在 Li 中不存在低于价带的 $p$ 带（由于 $p_x$ 是矢量，所有重要的矩阵元都是到 $p$ 带的）；但一般而言未必如此，而且事实上在其余碱金属中——可能 Na 除外——情形都不是如此。因此，到那些代表闭合壳层的带的矩阵元就相当重要。

虽然 Wigner 和 Seitz 的原始计算直接使用了 $k\cdot p$ 理论，但 Bardeen (27) 很快就引入了一个简单得多也精确得多的程序，他注意到波函数的一级修正可以基本上精确地算出。做法是直接迭代波动方程，即从展开
\begin{equation*}
u_k=u_0(r)+u_1^{\,k}(r)+\cdots
\end{equation*}
出发，于是有
\begin{equation*}
\left(-\frac{\hbar^2}{2m}\nabla^2+V-E_0\right)u_1=\frac{i\hbar^2}{m}\,k\cdot\nabla u_0
\end{equation*}
他观察到，$u_1$ 的一个特解是
\begin{equation*}
u_1^{(a)}=-i\,k\cdot r\,u_0(r)
\end{equation*}
（代入即可验证；非齐次项与 $\nabla(k\cdot r)\cdot\nabla u_0$ 项正好相消）。但它并不满足边界条件；这个边界条件显然是（因为 $u_1$ 应当具有 $p$ 函数的对称性并且是周期的）：解在元胞边界上为零。

我们可以在该解上加上齐次方程的一个 $p$ 函数解，使之在元胞边界上为零：
\begin{equation*}
V_1=\frac{i\,k\cdot r}{r^2}\,P(r)
\end{equation*}
其中 $P$ 是径向波动方程在固定能量 $E_0$ 下的一个解：
\begin{equation*}
\frac{d^2 P}{dr^2}-\frac{2}{r^2}\,P(r)+\frac{2m}{\hbar^2}\,(E_0-V)\,P=0
\end{equation*}
于是
\begin{equation*}
u_1=i\,k\cdot r\left(P/r^2-u_0(r)\right)
\end{equation*}
这样
\begin{equation*}
\left.\frac{P}{r^2}\right|_{R_c}=\left.u_0\right|_{R_c}
\end{equation*}

我们立刻认识到，量子亏损方法（quantum-defect method）的设置正是为了直接给我们——通过从 $p$ 态的光学光谱作内插——波函数 $P(r)$ 的一切有趣的性质，至少在元胞的外部部分是如此。通过运用一系列巧妙的变换，Bardeen 得以把二级能量修正完全用波函数的这个一级修正表示出来，甚至还用它在元胞边界处的行为表示出来：
\begin{equation*}
\alpha=\frac{m}{m^*}=\gamma\left[\frac{r}{P}\,\frac{dP}{dr}-1\right]_{R_c}
\end{equation*}
而
\begin{equation*}
\gamma=-\frac{1}{3}\,R_c\,\bigl[u_0(R_c)\bigr]^2\Bigg/\left[\gamma\,u_0(r)\,\frac{d}{dr}\left(\frac{1}{r}\,\frac{\partial u_0}{\partial E}\right)\right]_{R_0,\,E_0}
\end{equation*}
（译注：此式分母中的 $\gamma$ 与求值下标 $R_0$ 均按原书排印照录。）这个证明（部分归功于 Silverman 等人）相当冗长，所以我不打算在这里讨论。

在所有碱金属中，微扰论都收敛得相当好，因为 $u_0$ 相当平滑（微扰即 $\alpha(k\cdot\nabla)u_0$）。在 Na 中，$u_0$ 在元胞 90\% 的范围内几乎是个常数。于是，动能修正亦即费米修正 $E_F$ 恰好就是自由电子气的那个值，再对修改后的有效质量加以改正。它就是
\begin{equation*}
\int_0^{k_f} k^2\,dk\,\frac{\hbar^2 k^2}{2m^*}\Bigg/\int_0^{k_f} k^2\,dk=\frac{3}{10}\,\frac{\hbar^2 k_F^{\,2}}{m^*}
\end{equation*}
其中 $k_F$ 由下式确定：
\begin{equation*}
n=\frac{k_F^{\,3}}{3\pi^2}
\end{equation*}
这可以写成
\begin{equation*}
E_F=\frac{2.21}{r_s^{\,2}}\left(\frac{m}{m^*}\right)\quad\mathrm{ryd}
\end{equation*}
其中
\begin{equation*}
\frac{4\pi}{3}\,a_0^{\,3}\,r_s^{\,3}=\frac{1}{n}
\end{equation*}
对若干种碱金属给出各种相关的量（见表 3），会使我们对各种效应的数量级有一个感受。

% 注：p067 末句完整结束；其后正文所引“表 3”的表体在后续页，由下一文件处理。

% ---- B.2 收尾：表 3 与讨论段（原书 p.53 / PDF p068 顶部，协调者补译） ----
\begin{table}[H]
\centering
{\bfseries 表~3}\\[6pt]
\small
\setlength{\tabcolsep}{4.5pt}
\renewcommand{\arraystretch}{1.25}
\begin{tabular}{@{}llllllll@{}}
\toprule
元素 & $r_s$ & $-E_0$ & \shortstack{内聚边界修正\\$E_I-E_0$\ensuremath{,} ev}
 & \shortstack{动能修正\\$m/m^*$} & $E_F$\ensuremath{,} ev
 & \shortstack{未修正内聚能\\$E_I-E_0-E_F$\ensuremath{,} ev}
 & \shortstack{观测内聚能\\ev}\\
\midrule
Li & 3.21 & \shortstack[r]{0.69 ryd =\\9.4 ev} & 4.0 & 0.67 & 1.90 & 2.10 & 1.65\\
Na & 3.96 & \shortstack[r]{0.625 ryd =\\8.5 ev} & 3.15 & 1.02 & 2.0 & 1.15 & 1.20\\
Rb & 5.20 & 6.27 ev & 1.65 & 1.10 & 1.23 & 0.42 & 0.85\\
\bottomrule
\end{tabular}
\end{table}

这张表足以显示各种结果的大致面貌。正如我们所猜想的，Fermi 动能修正大致达到边界修正的二分之一到三分之二，而二者的差按数量级的程度给出金属的结合能。注意，内聚能是问题中最小的能量——它是作为大得多的数字之差被算出来的——这几乎总是如此，也正是结合能计算中困难的主要根源。把能带算准到 0.1 ryd 左右并不难，这对了解能带结构的一般面貌通常已令人满意，因为能带结构的尺度往往是 rydberg 量级，至少也是 0.4 到 0.5 ryd；但结合能却是 0.1 或更小。
